\documentclass[11pt, xcolor=dvipsnames,svgnames,x11names]{article}
\usepackage{xcolor}
\usepackage{xspace}
\usepackage{url}
\usepackage[colorlinks=true,bookmarks=false,linkcolor=black,urlcolor=blue,citecolor=black]{hyperref}
\usepackage{fancyhdr}
\usepackage[top=1in,bottom=1.2in,left=1in,right=1in]{geometry}
\usepackage{multicol}
\usepackage{pgfplots}
\usepackage{tikz}
\usetikzlibrary{circuits.logic.US, patterns}
\usepackage{mathpazo}
\usepackage{biolinum} 
\usepackage[all]{tcolorbox}
\usepackage{../Lectures/rahul_math}
\renewcommand{\familydefault}{\sfdefault}


% --- Custom Color Palette from Image ---
\definecolor{headerRed}{HTML}{A3403D} % Dark red from image top
\definecolor{patternBg}{HTML}{F2D9D7} % Light pinkish bg from image bottom
\definecolor{binaryText}{HTML}{D1A7A4} % Muted color for binary digits

\renewcommand{\texttt}[1]{{\ttfamily\color{headerRed}#1}}


% --- Background Pattern Setup ---
\usepackage[pages=all]{background}
\usepackage{bm}
\backgroundsetup{
scale=1,
color=black,
opacity=1,
angle=0,
contents={%
  \begin{tikzpicture}[remember picture,overlay]
    % Top Header Bar
    \fill[headerRed] (current page.north west) rectangle ([yshift=-1.2cm]current page.north east);
    % Bottom Binary Background
    \fill[patternBg] (current page.south west) rectangle ([yshift=3.5cm]current page.south east);
    % Decorative Line (from image)
    \draw[headerRed, line width=2pt] ([xshift=1in, yshift=-2.5cm]current page.north west) -- ([xshift=8in, yshift=-2.5cm]current page.north west);
    % Binary Text Overlay (Sample)
    \node[anchor=south, binaryText, font=\ttfamily\scriptsize, inner sep=0pt, yshift=1.5cm] at (current page.south) {
        \begin{tabular}{c@{\hspace{0.5em}}c@{\hspace{0.5em}}c@{\hspace{0.5em}}c@{\hspace{0.5em}}c@{\hspace{0.5em}}c@{\hspace{0.5em}}c}
        0 0 0 0 1 1 1 1 0 0 0 0 1 1 0 1 1 1 0 0 0 1 \\
        1 0 0 1 0 1 0 0 0 1 0 0 1 1 0 0 1 0 1 1 \\
        0 0 1 1 0 0 1 0 1 1 1 0 1 1 0 1 0 1 0 0 0 \\
        The University of Alabama in Huntsville
        \end{tabular}
    };
  \end{tikzpicture}
}
}

% --- Header Styling ---
\makeatletter
\renewcommand{\maketitle}{%
    \vspace*{0.2cm}
    \begin{center}
        \begin{tcolorbox}[
            enhanced,
            colback=white,
            colframe=headerRed,
            arc=0mm,
            title={\textcolor{white}{\bfseries \@title}},
            coltitle=white,
            fonttitle=\bfseries\large,
            attach boxed title to top left={xshift=10pt, yshift=-\tcboxedtitleheight/2},
            boxed title style={sharp corners, colback=headerRed, frame hidden}
        ]
        \large
        \vspace{10pt}
\noindent
\begin{tabular}{@{}p{3.5cm}p{7cm}r@{}}
\textbf{Student Name:} & \underline{\hspace{7cm}} & \textbf{Points:} 20 \\
\end{tabular}

\vspace{10pt}

\noindent
\begin{tabular}{@{}p{3.5cm}p{7cm}r@{}}
\textbf{A Number:} & \underline{\hspace{7cm}} & \textbf{Duration:} 30 Minutes \\
\end{tabular}
        \end{tcolorbox}
    \end{center}
    \vspace{10pt}
}
\makeatother

\newcommand{\hw}{Quiz 3}
\newcommand{\duedate}{April 02, 2026}

\pagestyle{fancy}
\fancyhf{}
\lhead{\small CPE 487/587, Spring 2026}
\rhead{\small \textit{\hw}}
\cfoot{\thepage}
\renewcommand{\headrulewidth}{0pt}

\def\checkmark{~~~ \tikz\fill[Green,scale=0.8](0,.35) -- (.25,0) -- (1,.7) -- (.25,.15) -- cycle;}

%\def\checkmark{~~~ }


\begin{document}

\title{\hw}
\maketitle

\begin{enumerate}

    \item Consider the following code for the Generator:

                \begin{minted}[
                    fontsize=\small,
                    breaklines,
                    baselinestretch=1.0
                ]{Python}
class Generator(nn.Module):
    def __init__(self, in=28, height=28, width=28):
        super(Generator, self).__init__()
        self.in = in
        self.height, self.width = height, width
        
        self.model = nn.Sequential(
            nn.Linear(in, 256),
            nn.LeakyReLU(0.2),
            nn.Linear(256, 512),
            nn.LeakyReLU(0.2),
            nn.Linear(512, 1024),
            nn.LeakyReLU(0.2),
            nn.Linear(1024, width*height),
            nn.Tanh()
        )
    
    def forward(self, z):
        img = self.model(z)
        return img.view(img.size(0), 1, 
                        self.height, self.width)
                \end{minted}

Based on the above code, what's the dimension of the noise that Generator uses in the default settings for mapping it to a possible image?
\begin{enumerate}
\item 256
\item 28 \checkmark
\item 512
\item 1024
\end{enumerate}

\item Consider the Bayesian framework:

    \begin{align*}
    p(\theta | \mathcal{D}) = \frac{p(\mathcal{D} | \theta)p(\theta)}{p(\mathcal{D})}
    \end{align*}
        
        with $p(\mathcal{D}) = \int p(\mathcal{D} |\theta)p(\theta)d\theta$.

        Here, $\theta$ is a parameter or latent representation. $\Dcal$ is dataset.        
        
        Out of all distribuions, which is the most difficult to model?

\begin{enumerate}
\item $p(\theta | \Dcal)$
\item $p(\Dcal)$ \checkmark
\item $p(\theta)$
\item $p(\Dcal | \theta)$


\end{enumerate}

\item How many Neural networks are present in Vanilla GAN:
\begin{enumerate}
\item 3
\item 1
\item 2 \checkmark
\item 4
\end{enumerate}




\item Explain cycle-consistency as you learnt in Cycle GAN.

\begin{tcolorbox}[colback=white, colframe=headerRed!75!black]
\vspace{1.5em}
\noindent\rule{\textwidth}{0.4pt}\\[1.5em]
\noindent\rule{\textwidth}{0.4pt}\\[1.5em]
\noindent\rule{\textwidth}{0.4pt}\\[1.5em]
\noindent\rule{\textwidth}{0.4pt}
\vspace{0.5em}
\end{tcolorbox}


\item Explain what are some issues encountered in GAN training.

\begin{tcolorbox}[colback=white, colframe=headerRed!75!black]
\vspace{1.5em}
\noindent\rule{\textwidth}{0.4pt}\\[1.5em]
\noindent\rule{\textwidth}{0.4pt}\\[1.5em]
\noindent\rule{\textwidth}{0.4pt}\\[1.5em]
\noindent\rule{\textwidth}{0.4pt}
\vspace{0.5em}
\end{tcolorbox}


\item If we have an input matrix as 

\begin{align*}
    \begin{bmatrix}
    3 & 5\\
    7 & 9
\end{bmatrix}
\end{align*}

and we are looking to apply fractionally-strided convolution with a stride of $2$, then how does the input changes before applying convolution using  a kernel $K$.

\begin{enumerate}
\item \begin{align*}
\begin{bmatrix}
    3 & 0 & 5 & 0\\
    0 & 0 & 0 & 0\\
    7 & 0 & 9 & 0\\
    0 & 0 & 0 & 0\\
    \end{bmatrix}  \checkmark
\end{align*}
\item \begin{align*}
\begin{bmatrix}
    3 & 0 & 5\\
    0 & 0 & 0\\
    7 & 0 & 9
    \end{bmatrix}
\end{align*}
\item \begin{align*}
\begin{bmatrix}
    3 & 5 & 7 & 9
    \end{bmatrix}
\end{align*}
\item  \begin{align*}
\begin{bmatrix}
    3 \\ 5 \\ 7 \\ 9
    \end{bmatrix}
\end{align*}
\end{enumerate}


\item What does AutoEncoder do?

\begin{enumerate}
\item It creates a compressed representation of dataset. \checkmark
\item It adds noise to dataset.
\item It can generate sythentic samples never seen  before based on a training set.
\item It enhances image resolution.
\end{enumerate}

\item Which one is the correct formula for KL Divergence \textbf{\underline{from}} distribution $q$ \textbf{\underline{to}} distribution $p$?

\begin{enumerate}
\item \begin{align*}
KL( q(\theta) || p(\theta | \Dcal) ) = {\displaystyle \LARGE \int} p(\theta) \log \frac{p(\theta | \Dcal)}{q(\theta)} d\theta 
\end{align*}
\item \begin{align*}
KL( q(\theta) || p(\theta | \Dcal) ) = {\displaystyle \LARGE \int} p(\theta) \log \frac{q(\theta)}{p(\theta | \Dcal)} d\theta 
\end{align*}

\item \begin{align*}
KL( q(\theta) || p(\theta | \Dcal) ) = {\displaystyle \LARGE \int} q(\theta) \log \frac{q(\theta)}{p(\theta | \Dcal)} d\theta  \quad \checkmark
\end{align*}
\item \begin{align*}
KL( q(\theta) || p(\theta | \Dcal) ) = {\displaystyle \LARGE \int} p(\theta) \log p(\theta | \Dcal)
\end{align*}
\end{enumerate}

\item Mean-field variational family assumes that 
\begin{enumerate}
\item latent variables are independent of each other in the approximate posterior. \checkmark
\item latent variables are orthogonal to each other in the approximate posteriror.
\item latent variables are linearly dependent on each oher in the approximate posterior.
\item latent variables have stochastic.
\end{enumerate}

\item What do you understand by prior in Bayesian framework?
\begin{tcolorbox}[colback=white, colframe=headerRed!75!black]
\vspace{1.5em}
\noindent\rule{\textwidth}{0.4pt}\\[1.5em]
\noindent\rule{\textwidth}{0.4pt}
\vspace{0.5em}
\end{tcolorbox}

\item The ELBO is equal to the negative KL divergence up to an additive constant.
\begin{enumerate}
\item True \checkmark
\item False
\end{enumerate}

\item In the context of Variational Autoencoder which one is not true, assuming $x$ as the input, $z$ as the latent variable, and $q$ as tractable approximation of the original distribution $p$?
\begin{enumerate}
\item $q_w(z|x)$ maps $x$ to a probability distribution over a latent space.
\item $p(z|x)$ is the intractable posterior, that is, it is very difficult to model.
\item $p(x)$ is known and you can always write easily.  \checkmark
\item $q_w(z|x)$ is the probabilistic encoder.
\end{enumerate}

\item 

Consider a random variable
\begin{align*}
\varepsilon \sim \Ncal (0, \Ibf)
\end{align*}

then the new random variable

\begin{align*}
    \mathbf{z} = \bm\mu + \bm\sigma \odot \varepsilon,
\end{align*}

for constant $\bm\mu$ and $\bm\sigma$ is a 



\begin{enumerate}
\item A Gaussian distribution. \checkmark
\item A Uniform distribution
\item An intractable distribution
\item A direc-delta distribution
\end{enumerate}

\item 

Consider a random variable
\begin{align*}
\varepsilon \sim \Ncal (0, \Ibf)
\end{align*}

Imagine that $\xbf_0$ is the initial data distribution (the original data distribution). 

Then, in the forward diffusion process, we have 

\begin{align*}
    \xbf_{t+1}  = \sqrt{1 - \beta_t}\xbf_t + \sqrt{\beta_t}\varepsilon
\end{align*}

Here, the recursive relationship collaposes into a single step relative to $\xbf_0$:

\begin{align*}
\xbf_{t+1} = \sqrt{\bar{\alpha}_{t+1}}\xbf_0 + \sqrt{1 - \bar{\alpha}_{t+1}}\varepsilon
\end{align*}

Choose the correct option that provides relationship between $\alpha_i$ and $\beta_i$

\begin{enumerate}
\item $\alpha_i = \beta_i$, and $\bar{\alpha}_{t+1} = \sum_{i=0}^t \alpha_i$
\item $\alpha_i = 1 + \beta_i$, and $\bar{\alpha}_{t+1} = \sum_{i=0}^t \alpha_i$
\item $\alpha_i = 1 - \beta_i$, and $\bar{\alpha}_{t+1} = \prod_{i=0}^t \alpha_i$ \checkmark
\item $\alpha_i = 1 + \beta_i$, and $\bar{\alpha}_{t+1} = \prod_{i=0}^t \beta_i$
\end{enumerate}

\item Referring to the previous problem, suppose somebody told you that initial data $\xbf_0$ has distribution $\Ncal(m,s)$ where $s$ is the variance and $m$ is the mean.

In such a case, what is the mean/expectation of $\xbf_{t+1}$ in the forward diffusion process (Hint: use the property of expectations):

\begin{enumerate}
\item $\Ebb[\xbf_{t+1}] = \sqrt{\bar{\alpha}_{t+1}}m$ \checkmark
\item $\Ebb[\xbf_{t+1}] = \sqrt{{\alpha}_{t+1}}m$ 
\item $\Ebb[\xbf_{t+1}] = \sqrt{{\beta}_{t+1}}m$ 
\item $\Ebb[\xbf_{t+1}] = \sqrt{1- {\beta}_{t+1}}m$ 

\end{enumerate}

\item An autonomous equation doesn't depend explicitly  on 

\begin{enumerate}
\item Time variable $t$ \checkmark
\item Input variable $x$
\item Output variable $y$
\item Derivative of the output variable $\frac{dy}{dt}$.
\end{enumerate}

\item Which one is true for Brownian motion:

\begin{enumerate}
\item Any sample at time $t$ is Gaussian distributed.
\item Increment between two consecutive time is Gaussian distributed. \checkmark
\item Increments are dependent on the previous increment.
\item In Brownian motion, the variance of the increment decreases with time.
\end{enumerate}

\item When does SDE collapses to become ODE?

\begin{enumerate}
\item When the diffusion coefficient $\sigma_t = 0$. \checkmark
\item When the diffusion coefficient $\sigma_t = \sigma$, a constant.
\item When the derivative of diffusion coefficient $\frac{ d \sigma_t}{dt} = 0$.
\item When the derivative of diffusion coefficient $\sigma_t$ cannot be determined.

\end{enumerate}

\item In ODE interpretation of Diffusion Model, ODE is defined via a time-dependent vector field $u: [0,1] \times \Rbb^d \to R^d$. $u_t$ is the target vector field at time $t$. If we are to approximate $u_t$ using a neural network $u_t^\theta$, what loss function we are interested to train $u_t^\theta$?

\begin{enumerate}
\item $\Lcal_{FM}(\theta) = \Ebb_{t, X_t}\bigg[ || u_t^\theta(X_t) - u_t(X_t)||^2 \bigg]$ \checkmark
\item $\Lcal_{FM}(\theta) = \Ebb_{t, X_t}\bigg[ || u_t^\theta(X_t) + u_t(X_t)||^2 \bigg]$
\item $\Lcal_{FM}(\theta) = \Ebb_{t, X_t} \bigg[|| u_t^\theta(X_t|X_{t-1}) - u_t(X_t)||^2 \bigg]$
\item $\Lcal_{FM}(\theta) = \Ebb_{t, X_t} \bigg[|| u_t^\theta(X_t|X_{t+1}) - u_t(X_t)||^2 \bigg]$
\end{enumerate}

\end{enumerate}

\end{document}